\documentclass[11pt]{article}
\usepackage[margin=1in]{geometry}
\usepackage{amsmath,amssymb,amsthm}
\usepackage{booktabs}
\usepackage{hyperref}

\newtheorem{lemma}{Lemma}
\newtheorem{proposition}[lemma]{Proposition}
\theoremstyle{remark}
\newtheorem{remark}[lemma]{Remark}

\newcommand{\cov}{\mathrm{ct}}
\newcommand{\cnt}{\mathrm{cnt}}

\title{Local search and a construction scan for\\covering-design and lottery-covering tables}
\author{Wynand van Staden\\{\small Independent researcher, South Africa}\\{\small\texttt{vanstadenw24@gmail.com}}\\{\small ORCID: \href{https://orcid.org/0009-0009-9884-8547}{0009-0009-9884-8547}}}
\date{27 October 2026}

\begin{document}
\maketitle

\begin{abstract}
A covering design $C(v,k,t)$ is a family of $k$-subsets (blocks) of a $v$-set containing every $t$-subset; a lottery
covering $C(v,k,t,m)$ only requires every $m$-subset to meet some block in at least $t$ points. We describe the
methods behind about a thousand improvements of the public tables of both. For covering designs, a decremental tabu
search with clause weighting, as in set-cover local search, keeps exact move scores incrementally. For lottery
coverings, the same search scores moves at the level of $t$-subsets by aggregating $m$-subset weights, so that a move
is scored with as many lookups as for covering designs (updating after a move remains costly); we also describe engineering for covers with up to
$2\cdot 10^9$ $m$-subsets. Finally, a dynamic-programming scan applies standard recursive constructions to the whole
live table, and the explicit products are then reduced by the search. We evaluated the searches in a test whose
protocol and sample were fixed in advance: 300 random covering-design entries and 300 random lottery entries not held
by us, two minutes per run. The weighted search improved 14\% of the covering-design entries and the same search
without weighting 8\%. On lottery entries the methods improved 1.3--4.3\%,
and the $t$-level search did better without weighting than with it. On 25 October 2026 the Covering Repository held
998 of our coverings: 401 covering designs and 65 lottery coverings found by search, and 532 constructed coverings.
\end{abstract}

\section{Introduction}

Let $V$ be a set of $v$ points. A \emph{$(v,k,t,m)$ lottery covering} is a family $\mathcal{B}$ of $k$-subsets of
$V$, the \emph{blocks}, such that every $m$-subset $M\subseteq V$ satisfies $|B\cap M|\ge t$ for some
$B\in\mathcal{B}$. The smallest size of such a family is $C(v,k,t,m)$; we also write $C(v,k,t,m)$ for a covering. The
lottery reading is that if the $m$ drawn numbers lie among the $v$ chosen ones, some ticket matches at least $t$ of
them. These objects are also called lotto designs, with notation $L(n,k,p,t)$ for $n=v$, $p=m$
\cite{batevanrees1998,livanrees2002,li1999}, and general covering designs \cite{montecalvo2012}. For $m=t$ the
condition says that every $t$-subset lies in a block, and $C(v,k,t)=C(v,k,t,t)$ is the covering number of a covering
design \cite{millsmullin1992,gordonstinson2007}, bounded below by the Sch\"onheim bound
$\lceil\frac{v}{k}\lceil\frac{v-1}{k-1}\cdots\lceil\frac{v-t+1}{k-t+1}\rceil\cdots\rceil\rceil$ \cite{schonheim1964}.
Finding $C(v,k,t)$ is a unicost set-cover problem with one element per $t$-subset and one set per $k$-subset.

Since exact values are known only in small cases, tables record the smallest known coverings. The La Jolla Covering
Repository recorded covering designs for $v<100$, $k\le 25$, $t\le 8$ until March 2026, and its data are archived
\cite{ljcr}. The Covering Repository \cite{cr} keeps covering designs and lottery coverings, accepts uploads and checks
each one with its own validator.

\paragraph{Related work.} Upper bounds come from algebraic and recursive constructions
\cite{gkp1995,montecalvo2012}, greedy methods \cite{gkp1995}, simulated annealing \cite{nurmela1993}, tabu search with
fast incremental evaluation \cite{fadlaoui2011}, general upper-bound constructions for $C(v,k,t,m)$
\cite{bertolo2004}, and exact or constraint-based methods for small cases
\cite{margot2003,cushing2024}; lotto-design tables combine many of these \cite{livanrees2002}. Our search borrows
from local search for satisfiability and set cover: clause weighting, which raises the weight of every unsatisfied
clause so that the search is pushed out of local minima, as in the breakout method \cite{morris1993} and dynamic local
search for SAT \cite{saps2002}; and row-weighting local search for unicost set cover, RWLS \cite{rwls2015} and NuSC
\cite{nusc2024}, which also search decrementally (find a cover, delete a set, repair).

\paragraph{Contributions.}
\begin{itemize}
\item A decremental, clause-weighted local search for covering designs that combines the neighbourhood of covering-design
search \cite{nurmela1993,fadlaoui2011} with the weighting and decremental scheme of set-cover local search
\cite{rwls2015}, with exact incremental scores (Section~\ref{sec:plain}).
\item A systematic scan that applies known constructions \cite{gkp1995,montecalvo2012} to the whole live table by
dynamic programming, including the searches' own fresh records, followed by explicit construction, certification and
reduction of the products by the search (Section~\ref{sec:constructions}).
\item An evaluation on a random sample of table entries, with a protocol fixed in advance, comparing the search with
weighting switched off and with simulated annealing (Section~\ref{sec:random}).
\item An extension to lottery coverings that scores moves by $t$-subset aggregates of $m$-subset weights, which are
the make and break scores of SAT local search on the monotone formula of Lemma~\ref{lem:twolevel}
(Section~\ref{sec:lottery}), and engineering for very large covers (Section~\ref{sec:big}). On lottery coverings the
evaluation is mostly negative: all methods improved few entries, and the weighting did not help the $t$-level search.
\end{itemize}
Section~\ref{sec:verify} describes verification, Section~\ref{sec:random} evaluates the searches on a random
sample of table entries, Section~\ref{sec:results} reports the published coverings by method, Section~\ref{sec:fail}
says where the methods do not work, and Section~\ref{sec:discussion} sums up. Code, record lists and
construction recipes are public \cite{repo}.

\section{Notation}

For a $t$-subset $T$ and a family $\mathcal{B}$, let $\cov(T)=|\{B\in\mathcal{B}:T\subseteq B\}|$. For an
$m$-subset $M$, let $\cnt(M)=|\{T\subseteq M:|T|=t,\ \cov(T)>0\}|$, the number of covered $t$-subsets of $M$.

\begin{lemma}\label{lem:twolevel}
An $m$-subset $M$ is covered (that is, $|B\cap M|\ge t$ for some block $B$) if and only if $\cnt(M)>0$.
\end{lemma}
\begin{proof}
If $|B\cap M|\ge t$, any $t$-subset $T$ of $B\cap M$ lies in $B$ and in $M$, so $\cnt(M)>0$. Conversely, if a
$t$-subset $T\subseteq M$ lies in a block $B$, then $|B\cap M|\ge t$.
\end{proof}

A \emph{move} replaces a point $x$ of a block $B$ by a point $y\notin B$. Only the $\binom{k-1}{t-1}$ $t$-subsets of
$B$ that contain $x$ can lose a block, and only the $t$-subsets of $B-x+y$ that contain $y$ can gain one. A $t$-subset
\emph{flips} when it changes between covered ($\cov>0$) and uncovered.

\section{Search}\label{sec:search}

\subsection{Covering designs}\label{sec:plain}

The search continues two lines of work. From local search for set cover it takes clause weighting, tabu and
decremental search as in RWLS \cite{rwls2015}: the weights of uncovered elements grow after each step, a tabu rule
prevents immediate reversals, and a cover is shrunk by deleting a set and repairing. From local search for covering
designs \cite{nurmela1993,fadlaoui2011} it takes the representation and the neighbourhood: a fixed number of blocks,
and moves that replace one point of a block. The combination is needed because the sets of the set-cover formulation,
all $\binom{v}{k}$ $k$-subsets, cannot be listed, so moves cannot add or remove listed sets as in RWLS.

The state is a family of $b$ blocks, a weight $w(T)\ge 1$ on every $t$-subset (initially $1$), and $\cov(T)$. The
cost is the total weight of uncovered $t$-subsets. One step is:
\begin{enumerate}
\item \emph{Focus.} Pick an uncovered $t$-subset $T^\ast$: the heaviest of 4 sampled with replacement, or with
probability $0.05$ a uniformly random one.
\item \emph{Candidates.} For every block $B$ with $|B\cap T^\ast|=t-1$, consider bringing in the missing point $y$ of
$T^\ast$ and removing some $x\in B\setminus T^\ast$. If no block qualifies, the step is skipped.
\item \emph{Choice.} Take the candidate with the best weighted gain (weight newly covered minus weight newly
uncovered), ties broken uniformly at random. A point removed from a block may not re-enter that block during the next $6$
to $9$ steps (tabu tenure $7$ plus a random $0$--$3$, counted from the current step); if every candidate is tabu, a random candidate is taken.
\item \emph{Weighting.} Add $1$ to $w(T)$ for every uncovered $T$.
\end{enumerate}
Step 4 applies the clause-weighting rule after every move rather than only at local minima. There is no restart and
no weight decay; a covered $t$-subset keeps its weight.

\paragraph{Decremental mode.} The search starts from a full covering, usually the current table entry. Whenever the
family is a covering it is saved, the block that covers the fewest $t$-subsets alone is deleted (ties to the lowest
index), the weights of the newly uncovered $t$-subsets are reset to $1$, and the search repairs. Runs stop at a time
limit or a target size.

\paragraph{Exact incremental scores.} Our first implementation scored each candidate by enumerating the affected
$t$-subsets; the second keeps the scores incrementally. For a block $B$ and $x\in B$, let $\mathrm{brk}(B,x)$ be the
total weight of the $t$-subsets $T$ with $x\in T\subseteq B$ and $\cov(T)=1$; this is the loss of removing $x$.
Weights change only on uncovered $t$-subsets, which never contribute to brk, so brk changes only when some $\cov(T)$
moves between $0$, $1$ and $2$: then $\mathrm{brk}(B',x)$ changes by $\pm w(T)$ for the sole owner $B'$ of $T$ and
each $x\in T$ (gain on $0\to 1$ and $2\to 1$, loss on $1\to 0$ and $1\to 2$). To find the owner in $O(1)$ we use a variant of the
critical-variable bookkeeping found in implementations of SAT local search (for example the code accompanying
probSAT \cite{probsat2012}), which track, for each clause with one true literal, the variable that makes it true; we
do not claim it as new. Give each
block a fixed identifier and let every $t$-subset store $\sigma(T)=\sum_{B\ni T}\mathrm{id}(B)$, updated whenever a
block gains or loses $T$.

\begin{lemma}
If $\cov(T)=1$, the unique block containing $T$ has identifier $\sigma(T)$.
\end{lemma}
\begin{proof}
The sum has one term.
\end{proof}

When a block is deleted, the last block takes its identifier and the $\sigma$ values of its $t$-subsets are corrected.
The gain of bringing $y$ into $B$ while removing $x$ is the weight of the uncovered $t$-subsets $S\cup\{y\}$ with $S$ a
$(t-1)$-subset of $B-x$. It is computed per candidate block, not stored: one pass over the $(t-1)$-subsets $S$ of $B$
accumulates the total $G$ and, for each point $p\in B$, the part $g(p)$ from sets $S\ni p$, so that the gain for
outgoing $x$ is $G-g(x)$. This costs $O(t\binom{k}{t-1}+k)$ for all $k$ outgoing points together, independently of
$v$. A build with consistency checks recomputes $\cov$, $\sigma$ and brk by brute force every 5{,}000 steps.

Table~\ref{tab:v1v2} compares the two implementations of the same search from the same start coverings. Incremental
scores give 11--18 times more steps per second and better coverings in the same time. The table also gives the
Sch\"onheim bound; the records are 22--71\% above it.

\begin{table}[t]
\centering
\begin{tabular}{lrrrrrr}
\toprule
 & & \multicolumn{2}{c}{blocks after 2 min} & \multicolumn{2}{c}{steps per second} & \\
\cmidrule(lr){3-4}\cmidrule(lr){5-6}
Entry & Start & recompute & incremental & recompute & incremental & Sch\"onheim \\
\midrule
$C(48,9,3)$ & 292 & 275 & \textbf{273} & 5{,}916 & 66{,}463 & 224 \\
$C(54,13,4)$ & 746 & 668 & \textbf{666} & 391 & 7{,}112 & 536 \\
$C(70,12,3)$ & 352 & 343 & \textbf{340} & 2{,}027 & 36{,}166 & 257 \\
$C(35,18,6)$ & 223 & 211 & \textbf{207} & 12 & 200 & 121 \\
$C(90,12,3)$ & 704 & 683 & \textbf{679} & 1{,}472 & 25{,}745 & 548 \\
\bottomrule
\end{tabular}
\caption{Scores recomputed per candidate versus kept incrementally; two-minute runs from the same coverings, one seed
each.}
\label{tab:v1v2}
\end{table}

\subsection{Lottery coverings}\label{sec:lottery}

The direct generalisation puts the weights on $m$-subsets and scores each move by enumerating the $m$-subsets it
affects. It is correct but slow: in our first tests, with 12 sampled candidate blocks, it removed a block from one of 78
record coverings. Lemma~\ref{lem:twolevel} allows scoring at the level of $t$-subsets instead, with 10--25 times more
steps per second (Appendix~\ref{app:tuning}). With the larger candidate sample of our batch runs, however, direct
scoring improved more random entries than the weighted $t$-level search (Section~\ref{sec:random}).

\paragraph{State.} We maintain $\cov(T)$ for every $t$-subset, and $\cnt(M)$ and a weight $w(M)$ for every
$m$-subset, together with the aggregates
\begin{align*}
A(T) &= \sum\{w(M) : T\subseteq M,\ \cnt(M)=0\}, &
L(T) &= \sum\{w(M) : T\subseteq M,\ \cnt(M)=1\},
\end{align*}
and $U(T)$, the number of uncovered $m$-subsets containing $T$. For uncovered $T$, $A(T)$ is the weight that covering
$T$ would gain. For covered $T$, an $m$-subset counted in $L(T)$ has $T$ as its only covered $t$-subset, so $L(T)$ is
the weight lost if $T$ becomes uncovered; $L$ is read only for covered $T$. A move $x\to y$ in block $B$ is scored with
$\binom{k-1}{t-1}$ terms each:
\[
\mathrm{loss} = \sum_{\substack{x\in T\subseteq B\\ \cov(T)=1}} L(T), \qquad
\mathrm{gain} = \sum_{\substack{y\in T\subseteq B-x+y\\ \cov(T)=0}} A(T).
\]

In SAT terms, Lemma~\ref{lem:twolevel} makes a lottery covering a monotone CNF formula with one Boolean variable per
$t$-subset (true when covered) and one clause per $m$-subset (the disjunction of its $t$-subsets); $\cnt(M)$ is the
number of true literals of a clause, and $A$ and $L$ are the weighted make and break scores of dynamic local search
for SAT \cite{saps2002}. What differs is that the variables are not set freely: a point swap flips a whole bundle of
$t$-subset variables at once, and the instance is far too large to write down. The search therefore sums per-variable
make and break scores over the bundle, which is where the estimate below comes from.

\begin{remark}\label{rem:approx}
This score is an estimate; the bookkeeping applied after a move is exact. It can err in three ways: (a) an
$m$-subset containing two newly covered $t$-subsets is counted twice in the gain; (b) an $m$-subset with
$\cnt(M)\ge 2$ whose covered $t$-subsets are all lost is missing from the loss; two lost $t$-subsets (both contain $x$)
can share $t-1$ points, so their union of $t+1\le m$ points can lie in one $m$-subset; (c) an $m$-subset with
$\cnt(M)=1$ that loses its covered $t$-subset but gains another is counted as lost. The score is exact when $m=t$ or
when at most one $t$-subset flips. We have not measured how often the errors change the chosen move.
\end{remark}

\paragraph{Updates.} When a $t$-subset flips, $\cnt(M)$ changes by one for each of its $\binom{v-t}{m-t}$ supersets
$M$, and only the changes $0\leftrightarrow 1$ and $1\leftrightarrow 2$ of $\cnt(M)$ affect the aggregates. When
$\cnt(M)$ goes $0\to 1$, $w(M)$ leaves $A$ and $U$ decreases on all $\binom{m}{t}$ $t$-subsets of $M$, and $w(M)$
joins $L$; $1\to 0$ is the reverse; $1\leftrightarrow 2$ removes or adds $w(M)$ in $L$. In the first implementation
$L$ is kept on all $\binom{m}{t}$ $t$-subsets of $M$ (only the covered one is ever read); Section~\ref{sec:big} keeps
it on the covered one only.

\paragraph{Weight bumps in $O(1)$.} The clause-weighting step adds the same amount to every uncovered $m$-subset. We
keep a global counter $E$ and, for an uncovered $M$, a base weight $w_0(M)$ with true weight $w_0(M)+E$. Then
$A(T)=A_0(T)+E\cdot U(T)$ with $A_0(T)=\sum w_0(M)$ over the same sets, and a bump adds the bump size to $E$. When $M$
becomes covered its weight is frozen at $w_0(M)+E$; when it becomes uncovered again its base is set to its weight minus
$E$.

\paragraph{The search.} The loop is that of Section~\ref{sec:plain} with these differences. The focus is an uncovered
$m$-subset $M^\ast$, the heaviest of 4 sampled. Candidates are blocks meeting $M^\ast$ in $t-1$ points, with any of the
$m-t+1$ points of $M^\ast$ outside the block as incoming point; 12 candidate blocks are sampled by default and up to
1{,}000 in our batch runs. A removed point is tabu for the next $4$ to $6$ steps, there is no random-focus noise, and a step
whose candidates are all tabu is skipped. If no block meets $M^\ast$ in $t-1$ points, a random block swaps a random point
outside $M^\ast$ for a random point of $M^\ast$, without tabu or weight bump. On each full covering all weights are
reset to $1$ and the block with the smallest $\sum_{T\subseteq B,\ \cov(T)=1}L(T)$ is deleted; this counts the
$m$-subsets whose only covered $t$-subset lies only in $B$, a lower bound on what $B$ covers alone.

A consistency-check build recomputes $\cov$, cnt, $A$ and $L$ by brute force every 2{,}000 steps. The implementation
needs about 14 bytes per $m$-subset during search (more at start-up) and refuses $\binom{v}{m}>3\cdot 10^8$. The
effects of the main settings (candidate sample, stopping rule, seeds and tenures) are listed in
Appendix~\ref{app:tuning}.

\section{Large covers}\label{sec:big}

Constructed covers can have $10^5$ blocks or more and $\binom{v}{m}$ near $10^9$. A second implementation keeps the
search of Section~\ref{sec:lottery} and changes the data structures; it accepts $\binom{v}{m}\le 2\cdot 10^9$ and also
runs covering designs ($m=t$). Its control settings (deletion pool, weight growth, reverts) are in
Appendix~\ref{app:big}.

\paragraph{Witness byte.} $L$ is now kept only on the covered $t$-subset of each $M$ with $\cnt(M)=1$, so a change
$1\leftrightarrow 2$ of $\cnt(M)$ must find that $t$-subset. Sort $M$, number the $t$-subsets $T\subseteq M$ by the
colex rank $\lambda(T)\in\{0,\dots,\binom{m}{t}-1\}$ of the positions of $M\setminus T$ in the sorted list, and keep
one byte $s(M)=\sum\lambda(T)\bmod 256$ over the covered $t$-subsets $T$ of $M$.

\begin{lemma}\label{lem:witness}
If $\binom{m}{t}\le 256$ and $\cnt(M)=1$, the covered $t$-subset of $M$ is the one with $\lambda=s(M)$.
\end{lemma}
\begin{proof}
The sum has one term $\lambda<\binom{m}{t}\le 256$, so reduction mod 256 leaves it unchanged, and $\lambda$ determines
$M\setminus T$ and hence $T$.
\end{proof}

A $1\leftrightarrow 2$ change then costs one update instead of $\binom{m}{t}$ (on $2\to 1$ the byte is updated before
it is read, on $1\to 2$ after). The program refuses $\binom{m}{t}>256$.\footnote{The check was added while writing this
paper. Without it, for $\binom{m}{t}>256$ the byte could name the wrong $t$-subset, so $L$ was kept on the wrong
$t$-subset and the bookkeeping, not only the scores, was wrong. Four published coverings were reduced by the unchecked
build: $C(39,9,7,11)$, $C(30,10,10,19)$, $C(31,10,10,20)$ and $C(32,10,10,20)$; the last three also exceed the flip-cost
limit below, which was raised for those runs. Like every covering the search
reduced, they were verified exhaustively afterwards (Section~\ref{sec:verify}), so they are valid.} It also refuses
parameters where one flip would cost more than $10^8$ updates, that is $\binom{v-t}{m-t}\binom{m}{t}>10^8$.

\paragraph{Parallel updates.} A flip enumerates the supersets of $T$ in parallel chunks. Each $m$-subset is visited by
exactly one thread, so per-$m$-subset state needs no synchronisation; shared $t$-subset aggregates use atomic additions
and the list of uncovered $m$-subsets a lock. Setup, score rebuilds, the candidate scan and candidate scoring are also
parallel above size thresholds; with several threads, runs are not reproducible from the seed.

Table~\ref{tab:big} shows the effect on one large cover.

\begin{table}[t]
\centering
\begin{tabular}{lrl}
\toprule
Version & steps/s & after 10 minutes \\
\midrule
first lottery implementation & -- & out of memory \\
large-cover version, parallel setup and deletion pool & $\approx 33$ & stalled near 366{,}000 \\
+ witness byte, reverts, growing bump & $\approx 120$ & \\
+ parallel flips, scoring and scan & 270--350 & $\approx 363{,}000$ \\
\bottomrule
\end{tabular}
\caption{$C(35,10,9,12)$ ($834{,}451{,}800$ $12$-subsets) from a constructed covering of $404{,}385$ blocks.}
\label{tab:big}
\end{table}

\section{A scan over standard constructions}\label{sec:constructions}

\subsection{Rules}

The following rules are standard; see \cite{gkp1995} for covering designs, \cite{montecalvo2012} for general coverings
and \cite{cushing2024} for disjoint unions of covering designs as lottery designs. We state them in the form the scan
uses. Blocks may be built with fewer than $k$ points; they are completed to $k$ points with arbitrary further points at
the end, which can only help. Correspondingly the value for $k>v$ is taken equal to that for $k=v$.

\begin{proposition}[Product on a split]\label{prop:product}
Let $V=V_1\sqcup V_2$ with $|V_1|=v_1$, $|V_2|=v_2$. Let $\mathcal{D}_i$ be a $(v_i,k_i,t_i,m_i)$ covering on $V_i$,
with $k_1+k_2=k$ and $t_1+t_2\ge t$. Then $\{B_1\cup B_2:B_i\in\mathcal{D}_i\}$ covers every $m$-subset $M$ with
$m_1\le |M\cap V_1|\le m-m_2$.
\end{proposition}
\begin{proof}
If $|M\cap V_1|\ge m_1$, an $m_1$-subset of $M\cap V_1$ meets some $B_1$ in at least $t_1$ points; likewise
$|M\cap V_2|\ge m_2$ gives $B_2$. The blocks are disjoint, so $|(B_1\cup B_2)\cap M|\ge t_1+t_2\ge t$.
\end{proof}

An $m$-subset has $a=|M\cap V_1|$ in $[\mathrm{lo},\mathrm{hi}]=[\max(0,m-v_2),\min(m,v_1)]$. Products whose intervals
$[m_1,m-m_2]$ together contain every feasible $a$ form a $(v,k,t,m)$ covering of size at most the sum of their sizes
(duplicate blocks are merged in the explicit covering). A factor with $t_i=0$ is a single empty block, so
$C(v,0,0,0)=1$. Restricting to $t_1+t_2=t$ loses nothing: a term with $t_1+t_2>t$ can lower $t_2$ (or $t_1$) and, by
the third rule of Proposition~\ref{prop:mono}, not grow.

\begin{proposition}[Complement duality]\label{prop:dual}
$|\bar B\cap\bar M| = v-k-m+|B\cap M|$. Hence the complements of a $(v,v-k,t',v-m)$ covering with
$t'=\max(0,v-k-m+t)$ form a $(v,k,t,m)$ covering of the same size.
\end{proposition}
\begin{proof}
$|\bar B\cap\bar M|=v-|B\cup M|=v-k-m+|B\cap M|$. If $v-k-m+t\ge 0$ the two conditions are equivalent. Otherwise
$|B\cap M|\ge k+m-v>t$ for every $B$ and $M$, and any single block is a covering.
\end{proof}

\begin{proposition}[Monotonicity]\label{prop:mono}
$C(v,k,t,m)$ is at most each of $C(v,k,t,m-1)$; $C(v,k-1,t,m)$; $C(v,k,t+1,m)$; $C(v+1,k,t,m)$; and
$C(v-1,k,t,m-1)$.
\end{proposition}
\begin{proof}
An $m$-set contains an $(m-1)$-set; a block can be completed; meeting $M$ in $t+1$ points implies meeting it in $t$.
Deleting the extra point from a covering on $v+1$ points (and completing the blocks) gives a covering on $v$ points,
since every $M\subseteq V$ meets each block in the same points. An $m$-subset of $V$ has at least $m-1$ points among
the first $v-1$.
\end{proof}
We call the last two rules (one point more, or one point fewer together with $m-1$) the \emph{one-point rules}.

\subsection{The scan}

The scan keeps a value $c(v,k,t,m)$ for $v\le 90$, $k\le 25$, $t\le m\le 20$. It starts from trivial values ($1$
if $t=0$ and $m\le v$; $0$ if $m>v$; $1$ if $v\le k$; $\lceil(v-m+1)/k\rceil$ if $t=1$) and from the smallest
covering available for each entry, in the live table or among our own coverings. Then, for up to six rounds and for
each $v$ in increasing order, it applies Proposition~\ref{prop:mono}, the best product (for $2\le t\le k<v$) and
complement duality. For the product, fix $v_1\le v/2$ and let
\[
W(m_1,m_2)=\min_{k_1+k_2=k,\ t_1+t_2=t} c(v_1,k_1,t_1,m_1)\,c(v_2,k_2,t_2,m_2),
\]
made non-increasing in both arguments by running minima (a term with smaller $m_i$ also serves larger ones). The
recursion
\[
F(\mathrm{lo})=0,\qquad F(r+1)=\min_{\mathrm{lo}\le l\le r}\bigl(F(l)+W(l,\,m-r)\bigr)
\]
finds the cheapest partition of $[\mathrm{lo},\mathrm{hi}]$ into consecutive intervals $[l,r]$, each covered by one
term; since $W$ is non-increasing, no cover by overlapping intervals is cheaper. The value for $v_1$ is
$F(\mathrm{hi}+1)$. The scan records the rule and $v_1$ used for each improvement; the construction tree is recovered
by recomputing the terms, expanded into explicit blocks, and certified (Section~\ref{sec:verify}).

\paragraph{Example.} $C(60,12,6,7)$ had $82{,}868$ blocks in the table. With $v_1=16$ the scan found six terms
(Table~\ref{tab:example}), $58{,}803$ blocks in total; the search then reduced the explicit covering to $58{,}228$.

\begin{table}[t]
\centering
\begin{tabular}{lllr}
\toprule
on $V_1$ ($v_1=16$) & on $V_2$ ($v_2=44$) & $a=|M\cap V_1|$ & blocks \\
\midrule
$C(16,12,6,7)=4$ & $C(44,0,0,0)=1$ & $7$ & 4 \\
$C(16,10,5,5)=34$ & $C(44,2,1,1)=22$ & $5$--$6$ & 748 \\
$C(16,8,4,4)=30$ & $C(44,4,2,3)=78$ & $4$ & 2{,}340 \\
$C(16,6,3,3)=38$ & $C(44,6,3,4)=387$ & $3$ & 14{,}706 \\
$C(16,4,2,2)=20$ & $C(44,8,4,5)=1{,}175$ & $2$ & 23{,}500 \\
$C(16,0,0,0)=1$ & $C(44,12,6,6)=17{,}505$ & $0$--$1$ & 17{,}505 \\
\midrule
 & & & 58{,}803 \\
\bottomrule
\end{tabular}
\caption{The product construction for $C(60,12,6,7)$. Each row covers the 7-subsets $M$ with $a$ in the given range.}
\label{tab:example}
\end{table}

\paragraph{Reducing the products.} Products cover many $m$-subsets several times, so the explicit covering is
reduced. Lottery products were shrunk by the large-cover search, or, where it refuses the parameters, only stripped of
\emph{redundant} blocks (blocks that are the only cover of no $m$-subset) by the checker of Section~\ref{sec:verify}.
Constructed covering designs were only stripped. To decide which published components are worth downloading, the
scan is also run as if every table entry were available to us.

\section{Verification}\label{sec:verify}

\paragraph{Covering designs found by search} were checked by an independent script that counts every $t$-subset.

\paragraph{Lottery coverings found by search} were checked exhaustively by a separate counting program, or for some
entries by a CP-SAT model whose infeasibility shows that no $m$-subset is uncovered.

\paragraph{Constructed coverings} are certified by construction. A checker rebuilds the block set from the recipe with
its own code and asserts the hypothesis of every step (Propositions~\ref{prop:product}--\ref{prop:mono}); the
components are verified separately as above. Every constructed covering that was shrunk or stripped was also checked
by a second exhaustive program. It rejects malformed input and removes duplicate blocks, then counts for every $m$-subset the
blocks meeting it in at least $t$ points, in a 16-bit counter per $m$-subset (2 bytes each, so up to about
$10^{10}$ $m$-subsets on our machine). A wrapped counter can only report a covered set as uncovered or keep a
redundant block, never accept a non-covering. It then deletes redundant blocks one at a time in file order,
re-checking each against the current counts; deletions only increase the number of $m$-subsets each remaining block covers
alone, so one pass suffices, and the result is irredundant (not necessarily minimum). According to our batch logs, 399
of the 630 constructed coverings we listed were checked this way at their published size. The other 231 (218 lottery
coverings, 95 of them with $\binom{v}{m}>2\cdot 10^9$, and 13 covering designs) rest on the
construction certificate and the Repository's validator.

Finally, the Covering Repository validates every upload with its own program before accepting it.

\section{A random sample of table entries}\label{sec:random}

The published coverings of Section~\ref{sec:results} are entries on which a method succeeded, chosen by us; they say
neither how often the methods succeed on an arbitrary entry nor whether other methods would have succeeded too. To
answer both we ran a comparison whose protocol was committed to the repository before any run
(\texttt{baselines/random\_sample/PROTOCOL.md}). It fixed the population, the sampling, the methods, the budget, the
outcomes and three hypotheses, but no statistical tests or decision rules; every test below is therefore
exploratory. All runs used one desktop computer (AMD Ryzen 5 5600X, 6 cores and 12 threads; 28\,GB of memory
available to the programs).

\paragraph{Population and sample.} From the Covering Repository tables of 24 October 2026, an entry $(v,k,t,m)$ is
eligible if $k<v$, the table value is between 10 and 2{,}000, $t\ge 3$ for covering designs and $t\ge 2$ for lottery
coverings, and $\binom{v}{t}\le 3\cdot 10^7$ (covering designs) or $\binom{v}{m}\le 3\cdot 10^7$ (lottery
coverings), so that every compared program can run it. Entries then held by one of our coverings are excluded, since
there the start covering would be our own search output: 9.8\% of the otherwise eligible covering-design entries and
1\% of the lottery entries. Of the 394 excluded covering-design entries 388 were held by search coverings (almost
all from the weighted search) and 6 by constructed ones; of the 99 lottery entries, 62 and 37. Entries the weighted
searches had tried without success remain in the population, so the exclusion biases the sample against the weighted
searches, by an amount we cannot measure. The population is the small and medium entries held by others; it says
nothing about entries with large $\binom{v}{m}$ (Section~\ref{sec:fail}). Entries were stratified by $t$ (3, 4, 5,
$\ge 6$) for covering designs and by $m-t$ (1, 2, 3, $\ge 4$) for lottery coverings, put in random order with a fixed
seed, and the first 75 per stratum with an available start covering were taken (none was skipped): 300 covering-design
entries and 300 lottery entries, drawn separately. The protocol first fixed 25 per stratum; it was amended to 75 while
the first runs were in progress, before any outcome had been looked at. During the runs a concurrent-write bug garbled
some result lines; the runner was patched and those runs were repeated from the start (repository history). Stratum
sizes in the population differ, so the pooled rates below are rates in the sample; weighting the strata by their
population sizes changes them by at most one percentage point.

\paragraph{Methods.} Each method started from the table's covering and searched decrementally for 120 seconds of wall
time, single-threaded, with three seeds; 11 runs shared the machine. For covering designs the start covering is the
archived La Jolla covering, which for 64 entries is larger than the current table value and must first be brought
down to it. The methods are ablations within our own framework, not reimplementations of other authors' programs. For
covering designs: the incremental search of Section~\ref{sec:plain} (\emph{weighted}; tenure 7, noise 0.05); the same
program with the weight increase switched off (\emph{tabu}); and simulated annealing over the same moves in the style
of \cite{nurmela1993} (\emph{annealing}), whose temperature schedule was tuned on a single entry, $C(57,12,3)$. For
lottery coverings: the search of Section~\ref{sec:lottery} with 1{,}000 candidate blocks and tenure 5
(\emph{weighted}); the same with the weight increase switched off (\emph{tabu}); and direct $m$-subset scoring with
the same settings (\emph{direct}), which is also weighted. An entry counts as improved by a method if its best
covering over the three seeds is smaller than the live table value. Each of the 69 improved covering designs and 25
improved lottery coverings (best per entry and method) was checked exhaustively and is valid. None of these coverings
is counted in Section~\ref{sec:results}.

\begin{table}[t]
\centering
\begin{tabular}{lrrrrr}
\toprule
Stratum & entries & weighted & tabu & annealing / direct & any \\
\midrule
\multicolumn{6}{l}{\emph{covering designs} (third method: annealing)} \\
$t=3$ & 75 & 14 & 10 & 0 & \\
$t=4$ & 75 & 15 & 5 & 0 & \\
$t=5$ & 75 & 7 & 6 & 0 & \\
$t\ge 6$ & 75 & 7 & 4 & 1 & \\
all & 300 & 43 & 25 & 1 & 54 \\
 & & 14.3\% & 8.3\% & 0.3\% & 18.0\% \\
 & & (10.8--18.8) & (5.7--12.0) & (0.1--1.9) & (14.1--22.7) \\
\midrule
\multicolumn{6}{l}{\emph{lottery coverings} (third method: direct)} \\
$m-t=1$ & 75 & 2 & 10 & 4 & \\
$m-t=2$ & 75 & 1 & 2 & 2 & \\
$m-t=3$ & 75 & 1 & 1 & 1 & \\
$m-t\ge 4$ & 75 & 0 & 0 & 1 & \\
all & 300 & 4 & 13 & 8 & 16 \\
 & & 1.3\% & 4.3\% & 2.7\% & 5.3\% \\
 & & (0.5--3.4) & (2.5--7.3) & (1.4--5.2) & (3.3--8.5) \\
\bottomrule
\end{tabular}
\caption{Entries of the random sample improved within $3\times 120$ seconds, by method and stratum, with 95\% Wilson
intervals for the sample rates. ``Any'' counts entries improved by at least one method.}
\label{tab:random}
\end{table}

\paragraph{Covering designs.} Table~\ref{tab:random} gives the improvement rates. The weighted search improved 43
entries, the tabu variant 25 and annealing one. On the 236 entries whose start covering equalled the table value the
counts are 38, 23 and 1. Of the entries improved by exactly one of weighted and tabu, 29 were weighted-only and 11
tabu-only (exact McNemar test $p=0.006$; $p=0.026$ after Holm correction over the six pairwise improvement tests of
this section; the sign tests below are not included in it, and none of them is significant). Compared on the median size over seeds, the weighted search was smaller than tabu on 36 entries, equal
on 239 and larger on 25 (sign test $p=0.20$), and smaller than annealing on 47 with no losses. Most entries are ties at
the table value, so the median relative gap is 0 for every method. Where the weighted search improved an entry, the
gain was small: median 0.71\%, mean 1.08\%. The advantage of weighting is not uniform. On entries with $k\ge 16$, 26
were weighted-only and 4 tabu-only, and on entries of 50--199 blocks 18 and 0; but on entries with $k<16$ it was 3
against 7, and on entries of 800 blocks or more 1 against 6. Annealing went below its start size in only 5 of 900
runs, so with this budget and schedule it serves only as a sanity check: the result says that our annealing settings
do not work from a table covering in two minutes, not that annealing cannot. The hypothesis stated in the protocol, that the
weighted search beats both alternatives on covering designs, holds against annealing, and against tabu on the
improvement rate overall but not on median size and not on every kind of entry.

\paragraph{Lottery coverings.} Every method improved few lottery entries, and most improvements came from $m-t=1$.
The $t$-level search improved 13 entries without weighting and 4 with it; 10 entries were tabu-only and 1
weighted-only (McNemar $p=0.012$, $p=0.035$ after Holm correction), and on median size tabu was smaller on 8 entries and
larger on 2 (sign test $p=0.11$). Direct scoring, which is weighted, improved 8 entries, between the two: 4 that the
weighted $t$-level search missed and none the other way ($p=0.12$), and against tabu it was smaller on median size on 3
entries and larger on 8 ($p=0.23$). So at these settings and this budget, weighting combined with the $t$-level score
estimates (Remark~\ref{rem:approx}) did worst; whether weighting also hurts direct scoring we did not test. The protocol
expected the $t$-level methods to beat direct scoring and weighting not to be clearly better than tabu. The second
hypothesis held, and the data point to weighting being worse. The first held for tabu, not significantly, and failed
for the weighted search. The improvements were mostly by one block (11 of 16 entries), the largest being
$C(45,12,4,5)$ from 226 to 210. Our lottery records in Section~\ref{sec:results} came mostly from the weighted search,
so some of them might also have been found without weighting.

\paragraph{Which entries yield.} The following counts for the weighted search on covering designs are descriptive; we
fitted no model. Success rose with $v$ (1 of 90 entries for $v<30$, 17 of 99 for $30\le v<50$, 11 of 71 for
$50\le v<70$, 14 of 40 for $v\ge 70$) and with $k$ (2 of 30 for $k<10$, 5 of 76 for $10\le k<16$, 36 of 194 for
$k\ge 16$). It also rose with the ratio of the table value to the Sch\"onheim bound (1 of 32 below 1.1, 23 of 121 at
1.6 or above). By table value it was low at both ends (1 of 91 below 50 blocks, 20 of 90 for 50--199, 20 of 72 for
200--799, 2 of 47 for 800 or more); at the top end two minutes is short for the weighted search, while the tabu variant
improved 7 of these 47. Where the start covering was the older La Jolla one and larger than the table value, 5 of 64
entries were improved, against 38 of 236 where it equalled the table value. On lottery coverings, where 16 entries were
improved by some method, only 6 of 144 entries with $v<30$ were improved, and 1 of 153 entries below 50 blocks. Small
entries, and lottery entries with small $v$, seem close to the best these methods reach.


\section{Published coverings}\label{sec:results}

This section reports what the methods produced on entries we chose ourselves, so it shows where they succeeded,
not how often. Table~\ref{tab:summary} summarises the published coverings on 25 October 2026. ``Listed'' counts every
covering we uploaded that improved its entry when uploaded; ``current'' counts those that still hold it. Of the 139
that no longer do, 135 (on 130 distinct entries) are now beaten by other contributors and 4 by our own constructed coverings. ``Our
value'' is our current best for the entry, which for 48 lottery and 9 search covering designs includes a later
reduction by the large-cover search.

\begin{table}[t]
\centering
\begin{tabular}{lrrrr}
\toprule
Source & listed & current & median saving & quartiles of saving \\
\midrule
covering designs, search & 432 & 401 & 1.24\% & 0.67--2.38\% \\
lottery coverings, search & 75 & 65 & 1.16\% & 0.43--2.36\% \\
constructed, product only & 357 & & 0.67\% & 0.29--1.70\% \\
constructed, product then reduced & 223 & & 1.48\% & 0.46--6.24\% \\
constructed, all & 630 & 532 & 0.85\% & 0.32--2.71\% \\
\midrule
total & 1{,}137 & 998 & & \\
\bottomrule
\end{tabular}
\caption{Published coverings on 25 October 2026. The saving of an entry is (table value before our upload $-$ our
value) / table value before our upload; quartiles by nearest rank. ``Reduced'' means shrunk by the search (128 lottery
coverings) or only stripped of redundant blocks (95 covering designs). Of the other 50 constructed coverings, 28 use
duality or a one-point rule at the top of their tree, and 22 are existing table coverings (ours or others') that were
shrunk or stripped.}
\label{tab:summary}
\end{table}

Table~\ref{tab:bymethod} breaks the same coverings down by the combination of methods that produced them.

\begin{table}[t]
\centering
\begin{tabular}{llrr}
\toprule
Kind & Method & listed & current \\
\midrule
covering designs & recompute search (first implementation) & 106 & 98 \\
 & \quad + later large-cover shrink & 4 & 4 \\
 & incremental search & 317 & 294 \\
 & \quad + later large-cover shrink & 5 & 5 \\
lottery coverings & lottery search & 27 & 20 \\
 & \quad + later large-cover shrink & 48 & 45 \\
\midrule
constructed covering designs & product & 150 & 147 \\
 & product, then redundant blocks removed & 95 & 94 \\
 & one-point rule, then redundant blocks removed & 5 & 5 \\
 & duality & 1 & 0 \\
 & existing covering, redundant blocks removed & 21 & 21 \\
constructed lottery coverings & product & 207 & 118 \\
 & product, then large-cover shrink & 128 & 124 \\
 & one-point rule, then large-cover shrink & 7 & 7 \\
 & duality & 15 & 15 \\
 & existing covering, large-cover shrink & 1 & 1 \\
\midrule
total & & 1{,}137 & 998 \\
\bottomrule
\end{tabular}
\caption{Published coverings by method on 25 October 2026. ``Current'' compares our value with the live table.}
\label{tab:bymethod}
\end{table}

\paragraph{Covering designs.} Of the 432 listed, 110 were first found by the recompute implementation and 322 by the
incremental one; 9 were later reduced further by the large-cover search. By $t$ they are $t=3$: 265, $t=4$: 112,
$t=5$: 38, $t=6$: 14, $t=7$: 3. From the batch logs: one pass of the incremental search over 3{,}251 table entries
improved its start covering on 319 of them, 312 of which improved the table (188 new records and 124 improvements of
our own earlier records); a second pass of 600-second runs from our own best coverings on those 319 shrank 249 further, by 1{,}629 blocks in total. The gains came
almost entirely from entries that other contributors had improved in 2025--26; the oldest entries (1996--2005) did not
improve.

\paragraph{Lottery coverings.} The 75 search records are for $(t,m)=(3,4)$ (39), $(4,5)$ (34) and $(4,6)$ (2); for
example $C(78,9,3,4)$ went from $672$ to $647$, and $C(74,8,3,4)$ from $784$ to $765$ and later, with the large-cover
search, to $756$. Among the search settings we tuned (Appendix~\ref{app:tuning}), none mattered as much as which entry was
attacked; the weighting itself, which we did not vary in these batches, mattered more (Section~\ref{sec:random}). The
following counts are from our batch logs, for entries we chose, and are anecdotal. On records set or re-optimised in 2025--26
the search succeeded on 0 of 102; on older ($\le 2023$) records by one prolific contributor or imported from earlier
literature with $t=4$, $m=5$, $v\ge 40$ it succeeded on 15 of 32. For $(t,m)=(3,4)$ from the same sources the rate was
roughly a tenth for $v\ge 40$ and near zero for $v<40$ (we did not record the denominators). This is the opposite of
what we saw for covering designs, where recently improved entries were the ones that yielded; we take it to depend on
how the competing entries were made rather than on their age, but did not test this.

\paragraph{Constructed coverings.} 580 of the 630 listed constructed coverings are products at the top of their tree,
and 223 of these were reduced afterwards. The gain from reducing varies widely: $C(56,14,7,8)$ fell from $143{,}599$
in the table to $106{,}701$ by a seven-term product and to $97{,}928$ after shrinking ($-8\%$), and $C(55,14,7,8)$ from
$122{,}870$ to $95{,}732$ and then $83{,}995$ ($-12\%$); but $C(60,6,5,7)$, built as $231{,}388\to 192{,}965$ by a
product with a one-point factor, shrank only to $192{,}940$.

\section{Where the methods do not work}\label{sec:fail}

Section~\ref{sec:random} measured how often the searches succeed on random entries of moderate size. What we observed
beyond it:
\begin{itemize}
\item \emph{Strong recent records.} In our batches (entries we chose; anecdotal) the lottery search improved none of
102 records set or re-optimised in 2025--26, and the covering-design search did not improve the oldest entries
(1996--2005).
\item \emph{Large $m$.} The search needs memory proportional to $\binom{v}{m}$, which rules out most lottery entries
(the median entry held by the table's largest contributor has about $7.6\cdot 10^{13}$ $m$-subsets). Those would need a
method that never enumerates $m$-subsets, such as lazy constraint generation with a fast exact checker. Exhaustive
verification has a similar limit, which is why 231 constructed coverings rest on their construction certificate.
\item \emph{Clause weighting with $t$-level estimates.} On the random sample the $t$-level search without weighting
improved more lottery entries than with it (13 against 4, Section~\ref{sec:random}), at one setting and one two-minute
budget. We have not tested other weighting rules, longer runs, or direct scoring without weighting.
\item \emph{Lottery entries in general.} The lottery methods improved 1.3--4.3\% of random entries, and only 6 of 225
entries with $m-t\ge 2$ were improved by any method.
\item \emph{Large covering designs in short runs.} On covering designs of 800 blocks or more the weighted search
improved 2 of 47 random entries in two minutes, fewer than the tabu variant (7).
\item \emph{Constructions.} Products rarely beat entries whose value comes from search, and reducing a product can
gain anything from almost nothing to 12\% (Section~\ref{sec:results}).
\item \emph{Ideas that gave nothing.} A two-for-one block exchange, restarting a level with another deletion, and
re-shrinking older coverings by other contributors (Appendix~\ref{app:tuning}).
\end{itemize}

\paragraph{Limitations.} All results are upper bounds; we prove no lower bounds and claim no optimality, and the gaps
to the Sch\"onheim bound remain large (Table~\ref{tab:v1v2}). Most timings are single runs, and several figures in
Section~\ref{sec:results} come from batch logs rather than from the published record lists. The lottery scores are
estimates (Remark~\ref{rem:approx}). The random sample covers only entries of moderate size not held by us, with one
short budget and one setting per method; longer runs or tuning could change the ranking, and the three methods share
our implementation of the framework, so the comparison is of scoring and weighting rules rather than of other
authors' programs. We have not measured how often the lottery score estimates change the chosen move.

\section{Discussion}\label{sec:discussion}

On covering designs, clause weighting in a decremental tabu search is the clearest positive result: on a random sample
of entries held by others it improved one entry in seven within two minutes, more than the same search without
weighting, mainly on entries with large $k$ and moderate size, and far more than our annealing settings. The gains are
small, typically under 1\%, and the success rate rises with the ratio of the table value to the Sch\"onheim bound;
small entries hardly moved.

The construction scan produced the largest number of records. It contains no new
construction: it applies known product, duality and monotonicity rules \cite{gkp1995,montecalvo2012} systematically to
a live table that contributors keep improving, and it pays off because a fresh record anywhere in the table can make
a product beat many other entries. Reducing the products by search then removes blocks the product covers several
times.

On lottery coverings the $t$-level scoring makes the search practical on large covers, but the random sample does not
support the clause weighting there: unweighted, the $t$-level search improved the most entries, though few in absolute
terms. Two directions follow. One is to understand why weighting combined with the score estimates of
Remark~\ref{rem:approx} does worse, by measuring how often the estimates mislead and by testing direct scoring
without weighting. The other is to reach the many lottery entries whose $\binom{v}{m}$ is far beyond enumeration, which
needs a search that never lists $m$-subsets.

\section*{Data and code availability}

\begin{sloppypar}
The code is at \cite{repo} under the MIT licence: \texttt{wls128} (recompute) and \texttt{wls2} (incremental) for
covering designs; \texttt{wlsm} (direct $m$-subset scoring), \texttt{wlsm2} (Section~\ref{sec:lottery}) and
\texttt{wlsm2big} (Section~\ref{sec:big}) for lottery coverings; the checkers \texttt{verify.py}, \texttt{mverify} and
\texttt{zerocov}; the scan \texttt{prodscan}; and \texttt{derived/build.py}, which rebuilds and checks any constructed
covering; and, in \texttt{baselines/random\_sample/}, the protocol, sampling script, sample, run script, raw results
and analysis script of Section~\ref{sec:random}. The repository also holds the record lists with previous values and the recipe of every constructed
covering. The component coverings are not redistributed, since coverings downloaded from the Covering Repository are
for personal use; each is identified by its parameters, size and SHA-1 hash and is available on the Covering
Repository, as are the coverings themselves.
\end{sloppypar}

\section*{Acknowledgements}

We thank Giovanni Acerbi for maintaining the Covering Repository and its validator, and the contributors whose
published coverings served as starting points and components. The code was written with the help of an AI assistant
(Claude, Anthropic); we designed the experiments, checked the results and are responsible for the content.

\appendix

\section{Tuning the lottery search}\label{app:tuning}

\begin{table}[h]
\centering
\begin{tabular}{p{7.2cm}p{6.6cm}}
\toprule
Change & Effect (our measurements) \\
\midrule
$t$-level scores and global weight counter & 10--25 times more steps per second than direct $m$-subset scoring \\
score up to 1{,}000 candidate blocks instead of 12 & about twice as many records found \\
stop an entry after 25\,s without shrinking & about 3.4 times more entries per hour, few records lost \\
three runs per entry (seeds, tabu tenure 2/5/10) & finds records single runs miss \\
two-for-one block exchange & correct, no extra records \\
restart a level with another deleted block & no measurable gain \\
\bottomrule
\end{tabular}
\caption{Effects measured while tuning the lottery search.}
\end{table}

\section{Control of the large-cover search}\label{app:big}

These are the settings of our batches; most are options, off by default.

\paragraph{Deletion.} The block to delete minimises the weighted cost $\sum_{T\subseteq B,\ \cov(T)=1}L(T)$. Candidates
come from a pool of the 2{,}000 blocks of lowest cost, rebuilt every 500 deletions; pool entries are re-scored in
order of their stored cost until the next stored cost is no better than the best found, or a cost of $0$ is found.

\paragraph{Weights.} Weights are kept between levels (by default they are reset to $1$), and the bump size $\delta$
grows by $\max(1,\lfloor 0.03\,\delta\rfloor)$ per level, which amounts to decaying older weights. When $\delta$
exceeds $2^{24}$, all weights and $\delta$ are reset.

\paragraph{Reverts.} If a deletion has not been repaired after 3{,}000 steps, or 120 seconds since the last full
covering or revert, or once the number of uncovered $m$-subsets exceeds $\max(1000,50u_0)$ where $u_0$ is the count
just after the deletion, the search returns to the last full covering through an undo log of moves. It then re-scores
all blocks with the current weights and deletes the block of rank $r\bmod 64$, where $r$ counts reverts so far.

\paragraph{Stopping and saving.} A run stops after 45--60 seconds without a smaller covering, or at 300--600 seconds.
Coverings are saved at most every 30 seconds and at exit, including on termination signals; our batch watchdog killed
overrunning runs and kept the last saved covering.

\begin{thebibliography}{99}

\bibitem{bertolo2004} R.~Bertolo, I.~Bluskov and H.~H\"am\"al\"ainen, Upper bounds on the general covering number
$C_\lambda(v,k,t,m)$, \emph{J. Combin. Des.} 12(5) (2004) 362--380.

\bibitem{batevanrees1998} J.~A. Bate and G.~H.~J. van Rees, Lotto designs, \emph{J. Combin. Math. Combin. Comput.}
28 (1998) 15--39.

\bibitem{cr} Covering Repository, \url{https://www.coveringrepository.com}, accessed 25 October 2026.

\bibitem{cushing2024} D.~Cushing and D.~I. Stewart, Applying constraint programming to minimal lottery designs,
\emph{Constraints} 29 (2024) 1--21. doi:10.1007/s10601-024-09368-5.

\bibitem{fadlaoui2011} K.~Fadlaoui and P.~Galinier, A tabu search algorithm for the covering design problem,
\emph{J. Heuristics} 17(6) (2011) 659--674. doi:10.1007/s10732-010-9150-2.

\bibitem{gkp1995} D.~M. Gordon, G.~Kuperberg and O.~Patashnik, New constructions for covering designs,
\emph{J. Combin. Des.} 3(4) (1995) 269--284. doi:10.1002/jcd.3180030404.

\bibitem{ljcr} D.~M. Gordon, La Jolla Coverings Repository, version 1.2, Zenodo, 2026.
doi:10.5281/zenodo.19735294. See also \url{https://dmgordon.org/covering-designs/}.

\bibitem{gordonstinson2007} D.~M. Gordon and D.~R. Stinson, Coverings, in C.~J. Colbourn and J.~H. Dinitz (eds.),
\emph{Handbook of Combinatorial Designs}, 2nd ed., Chapman \& Hall/CRC, 2007, Chapter VI.11.

\bibitem{li1999} B.~P.~C. Li, \emph{Some Results on Lotto Designs}, PhD thesis, University of Manitoba, 1999.
\url{http://hdl.handle.net/1993/1660}.

\bibitem{livanrees2002} P.~C. Li and G.~H.~J. van Rees, Lotto design tables, \emph{J. Combin. Des.} 10(5) (2002)
335--359. doi:10.1002/jcd.10020.

\bibitem{nusc2024} C.~Luo, W.~Xing, S.~Cai and C.~Hu, NuSC: an effective local search algorithm for solving the set
cover problem, \emph{IEEE Trans. Cybern.} 54(3) (2024) 1403--1416. doi:10.1109/TCYB.2022.3199147.

\bibitem{margot2003} F.~Margot, Small covering designs by branch-and-cut, \emph{Math. Program.} 94(2--3) (2003)
207--220.

\bibitem{millsmullin1992} W.~H. Mills and R.~C. Mullin, Coverings and packings, in J.~H. Dinitz and D.~R. Stinson
(eds.), \emph{Contemporary Design Theory}, Wiley, 1992, 371--399.

\bibitem{montecalvo2012} F.~Montecalvo, Some constructions of general covering designs, \emph{Electron. J. Combin.}
19(3) (2012) P28. arXiv:1205.4994.

\bibitem{morris1993} P.~Morris, The breakout method for escaping from local minima, in \emph{Proc. AAAI-93}, 1993,
40--45.

\bibitem{nurmela1993} K.~J. Nurmela and P.~R.~J. \"Osterg\aa rd, Upper bounds for covering designs by simulated
annealing, \emph{Congr. Numer.} 96 (1993) 93--111.

\bibitem{rwls2015} C.~Gao, X.~Yao, T.~Weise and J.~Li, An efficient local search heuristic with row weighting for the
unicost set covering problem, \emph{European J. Oper. Res.} 246(3) (2015) 750--761.
doi:10.1016/j.ejor.2015.05.038.

\bibitem{probsat2012} A.~Balint and U.~Sch\"oning, Choosing probability distributions for stochastic local search and
the role of make versus break, in \emph{Theory and Applications of Satisfiability Testing -- SAT 2012}, Lecture Notes
in Computer Science 7317, Springer, 2012, 16--29.

\bibitem{saps2002} F.~Hutter, D.~A.~D. Tompkins and H.~H. Hoos, Scaling and probabilistic smoothing: efficient dynamic
local search for SAT, in \emph{Principles and Practice of Constraint Programming (CP 2002)}, LNCS 2470, Springer,
2002, 233--248. doi:10.1007/3-540-46135-3\_16.

\bibitem{schonheim1964} J.~Sch\"onheim, On coverings, \emph{Pacific J. Math.} 14(4) (1964) 1405--1411.

\bibitem{repo} W.~van Staden, covering-design-records,
\url{https://github.com/WynandVStaden/covering-design-records}.

\end{thebibliography}

\end{document}
